% Experiments section. AAAI 2027 two-column, natbib.
% Style contract: no em dashes, no semicolon sentence joins, "we propose",
% every technical claim paired with a plain-English gloss.
% All numbers verified against results/*.json (see
% paper/figures/figure_data_provenance.md for the figure-side audit).

\section{Experiments}
\label{sec:experiments}

\subsection{Setup}
\label{sec:setup}

We evaluate on three headline tabular cells, named by their label-coupling
difficulty throughout: \textbf{easy} (HMDA, race $\to$ loan\_decision,
predictor $0.514$), \textbf{middle} (HMDA, race $\to$ loan\_amount\_band,
$0.584$), and \textbf{hard} (Adult, sex $\to$ income, $0.603$), plus two
CelebA vision cells selected by the same statistic. Attribute recovery is
measured as macro one-vs-rest AUC, where $0.5$ is chance and $0.55$ is the
protection bar. The recovery battery has two tiers, priced by attacker
knowledge. Tier~1 is black box: XGBoost \citep{chen2016xgboost}, a deep MLP
(256--256), and a rank-32 ReLU LoRA probe \citep{hu2022lora}. Tier~2 adds a
channel-aware Gaussian likelihood-ratio attacker that knows the clean
representation and the noise law. The LoRA member is not redundant: in the
end-to-end experiments it breached configurations that XGBoost and the MLP
both passed, so a battery without it would have certified operating points
that do not hold. Certified operating points are the mean over three
training seeds with the full battery on the representation and on the task
outputs, and every headline point is re-run at five seeds. In plain terms,
a method is only credited with hiding an attribute if the strongest probe
we can build, at the stated knowledge tier, reads near chance on every
exposed surface.

\subsection{Weak certificates leave a nonlinear gap}
\label{sec:demolition}

Figure~\ref{fig:demolition} is the audit of certificate-based erasure on
the Adult cell. Every targeted projection, whether it minimizes
class-conditional MMD in the style of LEOPARD or an HSIC independence
objective, drives the linear certificate to $R^2 \approx 0.036$, a clear
pass at $\tau = 0.05$. The honest probe reads the same representations at
$0.681$--$0.683$ AUC, and the unprotected representation reads $0.683$.
In plain terms, the projections removed what the certificate can see and
nothing that a gradient-boosted attacker uses. Only the noise channel
moves the honest number, from $0.589$ at $\sigma{=}1$ down to $0.514$ at
$\sigma{=}8$. On the HMDA cells the same projections never passed even the
certificate (attacked $R^2$ plateaus at $0.166$ with rank 32, against
$\tau = 0.05$), so Adult is the cell where the certificate-pass contrast
exists to be shown. Linear certificates, though deliberately scoped, leave
a large nonlinear gap that practical evaluations rarely check. We quantify
it: on this cell the gap is the full distance from a certified pass to
unprotected-level recovery.

\begin{figure}[t]
\centering
\includegraphics[width=\columnwidth]{figures/fig1_demolition.pdf}
\caption{Certificate-based erasure audited on Adult (sex). Left: the
linear certificate $R^2$, where every projection passes at
$\tau{=}0.05$. Right: honest recovery by XGBoost on the same
representations, where every projection sits at the unprotected level
($0.68$, top bar) and only the noise channel descends toward chance.
The certificate and the honest probe disagree about every projection and
agree only about noise. On HMDA the projections never passed even the
certificate, so no analogous contrast exists there.}
\label{fig:demolition}
\end{figure}

\subsection{Outputs leak what outputs actually use}
\label{sec:footprint}

Hiding the representation is not the whole problem, because the task
output is also an exposed surface. Figure~\ref{fig:footprint} plots, for
each cell, the best-attacker AUC on the task outputs once the
representation is hidden (the output-leak floor) against the
label-coupling predictor, a single XGBoost fit recovering the attribute
from the task label alone, computed before any model is trained. The six
tabular cells track the diagonal, from $(0.509, 0.509)$ to
$(0.676, 0.629)$. The causal law is that outputs leak what outputs
actually use. The predictor is therefore an upper bound on the output
floor, tight when the task is well learned and loose otherwise. The
CelebA Attractive$\to$Young cell makes the bound's slack measurable: the
label couples with Young at $0.734$, but the deployable model only
reaches $0.565$ accuracy, its achieved predictions couple with Young at
$0.536$--$0.540$, and the measured floor is $0.526$--$0.553$, matching
the achieved coupling rather than the label coupling. In plain terms, a
weak model cannot betray label structure it never captured. Where removal
is achievable the measured floors are small, $0.505$--$0.533$ on the
below-band cells ($0.526$--$0.553$ on the loose-bound CelebA cell),
empirically approaching the leakage floor established by
\citet{stadler2022}.

\begin{figure}[t]
\centering
\includegraphics[width=\columnwidth]{figures/fig2_footprint.pdf}
\caption{The output-leak floor against the label-coupling predictor, one
point per cell, with $y{=}x$. Tabular cells (Exp.~8) track the diagonal:
once the representation is hidden, the task output still betrays the
attribute to the extent the label couples with it. The CelebA cells
(Exp.~13) show the bound's character: the predictor is an upper bound on
the floor, tight when the task is well learned (Smiling$\to$Young, on the
diagonal) and loose when it is not (Attractive$\to$Young, far below, whisker
spanning the floors of the four erasure arms). There the achieved
predictions couple with the attribute at $0.536$--$0.540$ and the floor
matches that value, not the label coupling of $0.734$.}
\label{fig:footprint}
\end{figure}

\subsection{The cost of durable removal is a predictable ramp}
\label{sec:ramp}

Figure~\ref{fig:ramp} turns the predictor into a forecast of cost. For 20
natural cells across three datasets, the durable-removal cost, the
fraction of clean task lift lost at the noise level where the black-box
battery first reads at or below $0.55$ on both the representation and the
output, correlates with the predictor at $r = 0.795$ and
$\rho = 0.848$. Nine constructed cells, in which only the
label--attribute coupling is tuned, reproduce the same rise, which makes
the relationship causal rather than a dataset accident. The landscape is
a steep ramp, not a cliff: inside the band $0.52$--$0.60$ the cost climbs
from near zero to near everything, and outside it the answer is
essentially binary. In plain terms, one number computed from the labels
alone, before training anything, tells a practitioner whether durable
removal will be cheap, expensive, or hopeless.

\begin{figure}[t]
\centering
\includegraphics[width=\columnwidth]{figures/fig3_cost_ramp.pdf}
\caption{Durable-removal cost against the label-coupling predictor.
Filled circles are 20 natural cells ($r{=}0.79$, $\rho{=}0.85$), open
triangles are 9 constructed cells in which only the label--attribute
coupling is tuned, so the rise is causal. The shaded band marks the
usable middle ($0.52$--$0.60$), where the cost ramps from near zero to
near total. Below the band removal is nearly free, above it removal
costs essentially all task utility.}
\label{fig:ramp}
\end{figure}

\subsection{Two tiers: protection is priced by attacker knowledge}
\label{sec:twotier}

Figure~\ref{fig:twotier} prices certified protection at both tiers for
the end-to-end noised channel we propose. Tier~1 removes
representation-side leakage at near-zero cost on the achievable cells,
leaving only the output-level floor, which is fundamental: the easy cell
keeps $94.6\% \pm 2.2$ of its clean lift at $\sigma{=}8$ and the middle
cell keeps $56.6\% \pm 0.7$ at $\sigma{=}12$. Tier~2 is strictly more
expensive because the informed attacker must also be defeated: the same
cells certify at $\sigma{=}20$ and $\sigma{=}24$ for $55.5\% \pm 4.4$ and
$21.7\% \pm 0.5$. On the hard cell the diagnostic already predicts that
useful certification is out of reach, and the certified points confirm it
at $-0.5\%$ and $-7.0\%$. In plain terms, protection against a stronger
adversary is bought with more noise, and the price is read directly off
the utility axis.

Surgical destruction, which confines the noise to a learned
attribute-carrying subspace, is a Tier-1 technology. It holds the
black-box bar everywhere, but the informed likelihood-ratio attacker
reads its untouched complement at $0.66$--$0.94$, and no in-family
escalation fixes that. Only full-rank noise buys Tier~2. Within Tier~1
the surgical channel is what makes the hard cell usable at all: it keeps
$85\% \pm 2$ at the $0.55$ bar ($90\%$ at a relaxed $0.57$ bar), where
the $90\%$ point is the three-seed winner whose worst seed reads $0.561$
and breaches at five seeds. One scope note: the end-to-end advantage over
post-hoc surgery ($104\%$, $100\%$, and $90\%$ kept versus $88\%$,
$81\%$, and $53\%$ at matched Tier-1 protection, with the hard-cell
$90\%$ subject to the five-seed revision above) appears when
representation training is under our control. On the frozen, pre-shaped
CelebA representation the two tie ($99.1\%$ versus $98.4\%$ on the easy
vision cell).

\begin{figure}[t]
\centering
\includegraphics[width=\columnwidth]{figures/fig4_two_tier.pdf}
\caption{Utility kept at the certified operating point of the end-to-end
noised channel, per cell and per tier (error bars: spread over three
training seeds). Tier~1 (black box) is nearly free where the diagnostic
predicts removal is achievable, and what remains exposed afterward is
only the output-level floor, which is fundamental. Tier~2 (adding the
informed likelihood-ratio attacker) is bought only by more full-rank
noise, at $\sigma{=}20$/$24$/$64$, and its price is the visible drop. On
the hard cell the diagnostic predicts no useful certified point exists,
and none does.}
\label{fig:twotier}
\end{figure}

\subsection{The baseline gauntlet}
\label{sec:gauntlet}

We ran the field's published fair-representation and erasure methods
through the identical two-tier battery on the three headline cells, with
knob sweeps, best honest operating point per tier, and three-seed
certification. Each row carries a provenance tag (official code,
validated reimplementation, or generic mechanism), reported in the
appendix. One of 36 method--cell--tier combinations certifies.

\textbf{VFAE} \citep{louizos2016vfae} is the one survivor, and it
certifies through its accidental sampling noise rather than its fairness
machinery. The sampled-$z$ exposure holds Tier~1 on the easy cell at
$0.535$ with $72\%$ utility. The posterior-mean exposure of the same
models reads $1.00$ at every $\beta$, and raising the fairness knob makes
sampled-$z$ recovery strictly worse, $0.536$, $0.583$, $0.661$, $0.743$
for $\beta = 1, 10, 10^2, 10^3$, because a larger MMD weight shrinks the
posterior noise faster than it matches groups. Its Tier-2 reading of
$0.551$--$0.556$ is at the bar, not under it, and with no noise knob it
fails the middle and hard cells at $0.585$ and $0.587$. It is the
exception that proves the mechanism: the only baseline that certifies
anything does so exactly where the diagnostic says removal is cheap, via
an unadjustable noise channel it does not control.

\textbf{Obliviator} \citep{akbari2025obliviator}, run from the authors'
pinned official code in its strongest supervised mode, fails Tier~1 on
all three cells at every one of its 15 iterations, reading
$0.968$--$0.999$ while keeping $85$--$99\%$ utility, and its own MLP
stopping probe reads $0.94$--$0.99$ throughout, so its termination
criterion never fires. This matches the prediction of our controlled
HSIC experiment (Exp.~4): on representations where the attribute is
robustly and nonlinearly encoded, HSIC-guided erasure has almost nothing
kernel-visible left to remove, however seriously it is engineered.

\textbf{LAFTR} \citep{madras2018laftr}, both as our validated multiclass
reimplementation and re-run on the official Vector Institute release
under a compatibility shim calibrated to the paper's own Adult numbers,
fails Tier~1 everywhere, at $0.846$--$1.000$ across every adversary
weight. The official run adds one diagnosis: at $\gamma \ge 2$ on the
easy cell the min-max ``wins'' fairness by zeroing task utility, lift
falls to $3$--$17\%$ of clean, while the representation still leaks at
$0.996$--$1.000$. \textbf{DANN-style adversarial scrubbing}
\citep{ganin2015dann}, our generic-mechanism baseline at full strength,
genuinely removes signal but plateaus at $0.610$--$0.664$ for
$42$--$69\%$ utility. \textbf{LEACE} \citep{belrose2023leace}, from the
authors' library, reads $0.995$--$0.998$ at nearly full utility, the
unprotected level, as its linear scope predicts on representations that
encode the attribute nonlinearly. In plain terms, no published method
reaches the purpose-built channel on any cell at either tier, and the
single certifying combination is an accident of a stochastic encoder.
